\documentclass[a4paper]{article}
\usepackage{graphicx}
\usepackage{xcolor}
\usepackage{url}
\usepackage{outlines}
\usepackage{listings}
\usepackage{fontspec}
% \setmainfont{Calibri}
\lstset{basicstyle=\ttfamily,
	showstringspaces=false,
	commentstyle=\color{blue},
	keywordstyle=\color{pink}
}
\lstset{emph={
	nc,tcp,udp,http,},emphstyle=\color{purple}
}
\usepackage{fancyhdr}
\usepackage{geometry}
\geometry{
	a4paper,
	total={170mm,257mm},
	left=20mm,
	top=20mm,
	bottom=39mm,
}

\setlength{\headheight}{82.70538pt}

\fancypagestyle{oida}{
	\fancyhf{}
	\fancyhead[L]{\fontsize{7.5}{7.5}htl donaustadt\\ Donaustadtstraße 45\\
		1220 Wien\\~\\ Abteilung: Informationstechnologie\\ 
	Schwerpunkt: Netzwerktechnik}
	\fancyhead[R]{\includegraphics[scale=0.45]{images/logo.png}}

	\fancyfoot[L]{\today}
	\fancyfoot[C]{\jobname}
	\fancyfoot[R]{Seite: \thepage}
}

\begin{document}
\bibliographystyle{plain}
\pagestyle{oida}
\section*{Thema}
\par\noindent\rule{\textwidth}{0.4pt}

Laborprotokoll
Template

\begin{figure}[h]
	\includegraphics[scale=0.3]{images/mika.jpeg}
	\caption{Wunderbares Gruppenlogo}
\end{figure}

\vspace*{\fill}
Unterrichtsgegenstand:	NWT|ZIVK

Jahrgang:	3AHITN

Name:	Stefan Fürst, Marcel Raichle

Gruppenname/Nummer: Dumm und Dümmer/7

Betreuer: 	ZIVK

Übungsdaten:

Abgabedatum:	


\newpage
\tableofcontents

\newpage

\section{Aufgabenstellung}



\section{Zusammenfassung}

\newpage

\section{Vollständige Netzwerktopologie der gesamten Übung}

\newpage

\section{Übungsdurchführung}
%überschriften auf deutsch übersetzen
\subsection{Cable the Nework and Verify the Default Switch Configuration}
2b
	24 fe interfaces
		2 gbe ifs
		vty range
			5 15
2c 
	show startup-config
	startup-config is not present
	keine start conf wurde im NVRAM gespeichert
2d 
	keine ip in vlan1 per default
	mac addr
	0018.18bc.59c0
	its up
2e	
	show ip interface vlan1
	Internet protocal processing disabled
2f
	show ip interface vlan1
	Vlan1 is up, line protocol is up
	Internet protocol processing disabled
2g
	show version
	version:12.12
	image filename:
	show flash
	c3560-ipservicesk9-mz.122-55.SE6.bin
	mac addr:00:18:18:BC:59:80
2h
	show interface f0/6
	interface is up
	no shutdown/host connecten
	port mac 0018.18bc.5988
	Full-duplex, 100Mb/s,
2i
	Vlan1
	all of the ports bc every port is per default in vlan1
	its active
	Ethernet
2j 
	either show/dir flash
	show flash
		c3560-ipservicesk9-mz.122-55.SE6.bin
	dir flash

\subsection{Configure Basic Network Device Settings}
\paragraph {Configure basic switch settings}

a
no ip domain-lookup
hostname S1
service password-encryption
enable secret class
banner motd #
Unauthorized access is strictly prohibited. #

b
conf t
vlan 99 
exit
interface vlan99
ip address 192.168.1.2 255.255.255.0
ipv6 address 2001:db8:acad::2/64
ipv6 address fe80::2 link-local
no shutdown

c
interface range f0/1 – 24,g0/1 - 2
switchport access vlan 99

d
ss auf dc

e
ip default-gateway 192.168.1.1

f
line con 0 
logging synchronous
password cisco
login

g
ohne login command, wird nicht nach einem passwort gefragt


\paragraph {configure the pc ip addr}
step2:
ss in dc für beides
\subsection{Verify and Test Network Connectivity}
a
show run
b
show interface vlan99
bandwith: 1000000 Kbit
vlan99 state: up
line state: down

ping
https://discord.com/channels/@me/1052223143202279517/1296756855049289789


\subsection{Manage the MAC Address Tabel}

step1: 
D8-43-AE-86-06-88
step2:
show mac address-table
https://discord.com/channels/@me/1052223143202279517/1296757310357635115
wv dynamische mac addr
2
wv gesamt:
22
is dasselbe:
ja
step3:
a
13
b
1
show mac address-table dynamic
fixen und ss
c
show mac address-table address D8-43-AE-86-06-88
müssen fixen
step4:
a,b
clear mac address-table dynamic
show mac address-table
ss
c
1.
1 dynamic adresse
2. übesetzen sptäer: the switch dynamicly reacquired the pc mac address
d
mac address-table static d843.ae86.0688 vlan 99 interface fastethernet 0/6 
e
1.
22 mac addressen 
2. 
22 static
https://discord.com/channels/@me/1052223143202279517/1299273300035239946
f
no mac address-table static d843.ae86.0688 vlan 99 interface fastethernet 0/6
g
20 static oida weil wir eine removed haben :ReallyMad:
https://discord.com/channels/@me/1052223143202279517/1299274281036943393

reflection questions
Reflection Questions
1. Why should you configure the vty password for the switch?

If you do not configure a vty password you will not be able to telnet to the switch

2. Why change the default VLAN 1 to a different VLAN number?

For improved security

3. How can you prevent passwords from being sent in plain text?

Issue the service password-encryption command.

4. Why configure a static MAC address on a port interface?

To specify which ports a host can connect to.

https://discord.com/channels/@me/1052223143202279517/1296753070125879369
\newpage

\subsubsection{ssh}
Switch benötigt domain name da keys unter hostname.domain abgespeichert werden und eine hostnamen
- `ip domain-name htl-donaustadt.at
- `crypto key gen rsa
	- länge eingeben 4096 normalerweise aber im labor nur 1024 nehmen
- benutzer anlegen
	- `username name secret cisco
		- secret statt passwort so dass es keine oarsch verschlüsselung used
	- anmelden
		- `login local
			- locale benutzerdatenbank nutzen einschalten
		- `logging synchronous
			- nd required aba is qol
		- beim anmelde bildschirm fragt es dann nach nem pw
- `line vty 0 4` 4 weil 4 gleichzeitige verbindgen erlauben
	- `login local
	- `transport input ssh
		- ssh einschalten
	- `do wr
		- copy run start in besser

\subsubsection{aufräumen}

en
show flash
delete vlan.dat
enter drücken
erase startup-config
reload
no
\section{Vollständige Konfigurationsdateien}

Current configuration : 3608 bytes
!
version 12.2
no service pad
service timestamps debug datetime msec
service timestamps log datetime msec
service password-encryption
!
hostname S1
!
boot-start-marker
boot-end-marker
!
enable secret 5 $1$WG/H$ILJURiwvDuU5JlSgXmVwY/
!
username name secret 5 $1$RrXg$grWbyFauxAuHdmwpC69tK.
!
!
no aaa new-model
system mtu routing 1500
no ip domain-lookup
ip domain-name htl-donaustadt.at
!
!
!
!
crypto pki trustpoint TP-self-signed-662526720
 enrollment selfsigned
 subject-name cn=IOS-Self-Signed-Certificate-662526720
 revocation-check none
 rsakeypair TP-self-signed-662526720
!
!
crypto pki certificate chain TP-self-signed-662526720
 certificate self-signed 01
  30820239 308201A2 A0030201 02020101 300D0609 2A864886 F70D0101 04050030
  30312E30 2C060355 04031325 494F532D 53656C66 2D536967 6E65642D 43657274
  69666963 6174652D 36363235 32363732 30301E17 0D393330 33303130 30303130
  345A170D 32303031 30313030 30303030 5A303031 2E302C06 03550403 1325494F
  532D5365 6C662D53 69676E65 642D4365 72746966 69636174 652D3636 32353236
  37323030 819F300D 06092A86 4886F70D 01010105 0003818D 00308189 02818100
  E8CE9436 BED8F37A A6DCE351 5227D20F B07A6BC0 FA0445B9 CEDC0064 57A2B496
  39161F3F D82FAB21 BA4D34D7 4DB51AAF 0A42E5C1 93AC51A4 B61D11F8 9CC33A19
  51920ADB 3266102A 8A5745DA 06ABA47C FECC8C07 AF90612C 412CB8E3 F26E329C
  CC17F9E2 81D47732 B02C8AC9 33C82388 87D3E4DF 2E86B505 E4170470 021733ED

S1#sh run
Building configuration...

Current configuration : 3608 bytes
!
version 12.2
no service pad
service timestamps debug datetime msec
service timestamps log datetime msec
service password-encryption
!
hostname S1
!
boot-start-marker
boot-end-marker
!
enable secret 5 $1$WG/H$ILJURiwvDuU5JlSgXmVwY/
!
username name secret 5 $1$RrXg$grWbyFauxAuHdmwpC69tK.
!
!
no aaa new-model
system mtu routing 1500
no ip domain-lookup
ip domain-name htl-donaustadt.at
!
!
!
!
crypto pki trustpoint TP-self-signed-662526720
 enrollment selfsigned
 subject-name cn=IOS-Self-Signed-Certificate-662526720
 revocation-check none
 rsakeypair TP-self-signed-662526720
!
!
crypto pki certificate chain TP-self-signed-662526720
 certificate self-signed 01
  30820239 308201A2 A0030201 02020101 300D0609 2A864886 F70D0101 04050030
  30312E30 2C060355 04031325 494F532D 53656C66 2D536967 6E65642D 43657274
  69666963 6174652D 36363235 32363732 30301E17 0D393330 33303130 30303130
  345A170D 32303031 30313030 30303030 5A303031 2E302C06 03550403 1325494F
  532D5365 6C662D53 69676E65 642D4365 72746966 69636174 652D3636 32353236
  37323030 819F300D 06092A86 4886F70D 01010105 0003818D 00308189 02818100
  E8CE9436 BED8F37A A6DCE351 5227D20F B07A6BC0 FA0445B9 CEDC0064 57A2B496
  39161F3F D82FAB21 BA4D34D7 4DB51AAF 0A42E5C1 93AC51A4 B61D11F8 9CC33A19
  51920ADB 3266102A 8A5745DA 06ABA47C FECC8C07 AF90612C 412CB8E3 F26E329C
  CC17F9E2 81D47732 B02C8AC9 33C82388 87D3E4DF 2E86B505 E4170470 021733ED
  02030100 01A36330 61300F06 03551D13 0101FF04 05300301 01FF300E 0603551D
  11040730 05820353 312E301F 0603551D 23041830 16801460 47280E51 C2029C61
  DF3F8BB8 D9255894 39459730 1D060355 1D0E0416 04146047 280E51C2 029C61DF
  3F8BB8D9 25589439 4597300D 06092A86 4886F70D 01010405 00038181 00C0E85C
  A8F3A3E8 D613AA85 A036A6F4 3A1DD66B F05114A6 D03A2A06 620A9D2D 460E0F53
  B94F1B4F DE21ECDD BBD7C1D1 1A3C17B9 BAE76D91 C08AF26C FEBBBF0A 05A2653F
  596632F5 38D8C1C6 0C5A58F3 C90C797E 99E9E3AC 48AA92A4 F2F59711 6987BBA3
  A761D3BA AE41B89F 6933C814 87608BAD C087AB4C E681B888 73A5BE20 6E
  quit
!
!
!
spanning-tree mode pvst
spanning-tree extend system-id
!
vlan internal allocation policy ascending
!
!
!
!
interface FastEthernet0/1
 switchport access vlan 99
!
interface FastEthernet0/2
 switchport access vlan 99
!
interface FastEthernet0/3
 switchport access vlan 99
!
interface FastEthernet0/4
 switchport access vlan 99
!
interface FastEthernet0/5
 switchport access vlan 99
!
interface FastEthernet0/6
 switchport access vlan 99
!
interface FastEthernet0/7
 switchport access vlan 99
!
interface FastEthernet0/8
 switchport access vlan 99
!
interface FastEthernet0/9
 switchport access vlan 99
!
interface FastEthernet0/10
 switchport access vlan 99
!
interface FastEthernet0/11
 switchport access vlan 99
!
interface FastEthernet0/12
 switchport access vlan 99
!
interface FastEthernet0/13
 switchport access vlan 99
!
interface FastEthernet0/14
 switchport access vlan 99
!
interface FastEthernet0/15
 switchport access vlan 99
!
interface FastEthernet0/16
 switchport access vlan 99
!
interface FastEthernet0/17
 switchport access vlan 99
!
interface FastEthernet0/18
 switchport access vlan 99
!
interface FastEthernet0/19
!
interface FastEthernet0/20
!
interface FastEthernet0/21
!
interface FastEthernet0/22
!
interface FastEthernet0/23
!
interface FastEthernet0/24
!
interface GigabitEthernet0/1
!
interface GigabitEthernet0/2
!
interface Vlan1
 no ip address
 shutdown
!
interface Vlan99
 ip address 10.0.0.69 255.255.255.0
!
ip classless
ip http server
ip http secure-server
!
!
!
!
!
line con 0
line vty 0 4
 login local
 transport input ssh
line vty 5 15
 login
!
end
\newpage

\section{Quellen}
\bibliography{quellen}
\newpage
\section{Abbildungsverzeichnis}

\listoffigures

\end{document}
